\documentclass[10pt,a4paper]{amsart}
\usepackage[margin=19mm]{geometry}
\usepackage{amsmath,amssymb,mathtools}
\usepackage[scr=rsfs,cal=euler]{mathalfa}
\usepackage{xfrac}
\usepackage[unicode,hidelinks,bookmarks=false]{hyperref}
\newcommand{\ie}{i.e.\ }
\newcommand{\cf}{cf.\ }
\newcommand{\resp}{respectively\ }
\newcommand{\Subsection}{\subsection}
\providecommand{\nobreakdash}{\nobreak\hbox{-}}
\setlength{\emergencystretch}{2em}
\multlinegap0pt
\usepackage{amscd}
\newtheorem{theorem}{Theorem}
\newtheorem{lemma}[theorem]{Lemma}
\newtheorem{definition}[theorem]{Definition}
\newtheorem{corollary}[theorem]{Corollary}
\newtheorem{example}[theorem]{Example}
\newtheorem{remark}[theorem]{Remark}
\newcommand{\C}{\mathbb{C}}
\newcommand{\R}{\mathbb{R}}
\newcommand{\T}{\mathbb{T}}
\newcommand{\Pic}{\operatorname{Pic}}
\begin{document}
\title[A New CR Invariant for Contact 3-Manifolds and Classes of Op]{A New CR Invariant for Contact 3-Manifolds and Classes of Open Books}
\author{Ali M. Elgindi}
\date{}
\maketitle
\noindent\emph{Source and licence.} A New CR Invariant for Contact 3-Manifolds and Classes of Open Books. Version arXiv:2606.20750v2 (CC BY 4.0). This material is distributed under the Creative Commons Attribution 4.0 International licence, http://creativecommons.org/licenses/by/4.0/, as recorded in the arXiv OAI licence metadata for this exact version and shown on the abstract page https://arxiv.org/abs/2606.20750v2. Full rights notice: licenses/arxiv-2606.20750-CC-BY-4.0.txt.\par\smallskip
\noindent\textit{Editorial scope.} Counted unit: the original \texttt{theorem} environment \texttt{-} at source lines 97--99 together with its complete proof at lines 101--105; the delivered document keeps source lines 73--106 so that every referenced label and every citation resolves inside it. Native editable source is the paper's own concatenated TeX; the AMS wrapper is editorial formatting only. No publisher class, upstream script or unrelated artwork is redistributed.
\section{The Contact Complex Line Bundle}

Let $B$ be any supporting open book for $(M,\xi)$ and let $\gamma$ be its binding. By [1], there exists an embedding $M \hookrightarrow \C^3$ such that complex tangencies are exactly $\gamma$, and holomorphic tangent spaces along $\gamma$ are $\xi|_\gamma$. This induces a resulting complex line bundle $L_\xi$ over $M$ with respect to the choice of supporting book $B$:

\begin{definition}[The Picard Invariant]
For a contact structure $\xi$ on $M$, we define the Picard invariant as
\[
\mu_M(\xi) = [L_\xi] \in \Pic_{\C}(M).
\]
\end{definition}

\begin{lemma}[Stabilization Invariance]
\label{lem:stab}
Let $(B, \pi)$ be an open book supporting $\xi$, and let $(B', \pi')$ be a positive stabilization of $(B, \pi)$. Then the complex line bundles $L_\xi$ and $L'_\xi$ constructed from these open books via the embedding theorem \cite{Elgindi2025} are isomorphic.
\end{lemma}

\begin{proof}
Let $F: M \hookrightarrow \mathbb{C}^3$ be an embedding as constructed in [2]with respect to $(B,\pi)$ complex tangent along the link $B$. Let $F': M \hookrightarrow \mathbb{C}^3$ be another embedding given by our construction with respect to $(B',\pi')$ with complex tangent link $B'$. Both these embeddings exist by our paper \cite{Elgindi2025}.

Consider the holomorphic line bundles on $\mathbb{C}^3$ obtained by Stein extension of $\xi|_B$ and $\xi|_{B'}$ respectively. Call them $\mathcal{L}$ and $\mathcal{L}'$. Since $\mathbb{C}^3$ is Stein, Cartan's Theorem B implies that a holomorphic line bundle is uniquely determined by its first Chern class. The first Chern class $c_1(\mathcal{L})$ is determined by the CR structure of $\xi$ (via the contact form and the almost complex structure on $\xi$). Because the supporting structure $\xi$ is the same for both open books, we have $c_1(\mathcal{L}) = c_1(\mathcal{L}')$. Hence $\mathcal{L} \cong \mathcal{L}'$ as holomorphic line bundles over $\mathbb{C}^3$.

Restricting to $M$ gives an isomorphism $L_\xi \cong L_\xi'$. Therefore, the Picard invariant $\mu_M(\xi) = [L_\xi]$ is independent of the choice of open book.
\end{proof}


\begin{theorem}
The invariant $\mu_M(\xi)$ is well-defined and independent of the choice of supporting open book.
\end{theorem}

\begin{proof}
Let $\mathcal{B}_1, \mathcal{B}_2$ be open books supporting $\xi$ on $M$, with respective bindings $\gamma_1, \gamma_2$. Each book yields an embedding $f_i: M \hookrightarrow \C^3$ that is holomorphic in a neighborhood of $\gamma_i$ and gives a holomorphic line bundle $L_i$ as the holomorphic tangent spaces along $\gamma_i$. Since $\C^3$ is Stein, the line bundles $L_i$ extend to holomorphic line bundles over the ambient $\C^3$; restricting to $M$ gives holomorphic line bundles over $M$.

By Giroux's theorem, any two open books supporting the same contact structure $\xi$ are related by a sequence of positive stabilizations [7]. By Lemma \ref{lem:stab}, the line bundles obtained from each open book are isomorphic. Hence $\mu_M(\xi)$ is independent of the choice of supporting open book and gives a well-defined invariant of contact structures on closed 3-manifolds.
\end{proof}


\begin{thebibliography}{99}
\bibitem[Elgindi2025]{Elgindi2025} A. M. Elgindi, \textit{Explicit holomorphic structures for embeddings of 3-manifolds into $\C^3$}, Methods Appl. Anal. \textbf{32} (2025), no. 1, 1–21.
\end{thebibliography}
\end{document}
